\subsubsection{Retrievalprestatie en VLAM-poort van de advies-matcher.}
De advies-matcher zoekt adviesrapporten van Kaderwet-adviescolleges op in de
officiële bekendmakingen. Het systeem werkt in twee fasen. Een deterministische
retrievallaag haalt eerst kandidaat-documenten op door op de collegenaam te zoeken.
Daarna beoordeelt een VLAM-poort elke kandidaat inhoudelijk en bepaalt of het
werkelijk een adviesrapport is. De evaluatie hieronder kijkt naar beide fasen, met
de nadruk op de vraag of het juiste rapport \"uberhaupt in de kandidaatlijst
verschijnt.

\noindent Als grondwaarheid dienden adviesrapporten die rechtstreeks van de websites
van vijf grote, permanente colleges zijn binnengehaald: de Adviesraad Internationale
Vraagstukken, de Raad voor de Leefomgeving en Infrastructuur, de Onderwijsraad, de
Raad voor het Openbaar Bestuur en de Raad voor Cultuur. Deze bron is onafhankelijk
van de matcher, want de matcher zoekt uitsluitend in de officiële bekendmakingen.
Een match werd vastgesteld op tekstovereenkomst, namelijk de Jaccard-similariteit van
woord-shingles met een drempel van 0,40. Die keuze was nodig omdat hetzelfde document
in beide bronnen een ander document-ID heeft, maar vrijwel dezelfde tekst.

\noindent Niet elk website-rapport staat ook in de officiële bekendmakingen. Een
aanwezigheidsscan stelde vast dat 140 van de 403 website-rapporten daar aantoonbaar
aanwezig zijn. De recall is alleen over deze 140 berekend, zodat een gemist rapport
een echte fout van de matcher is en geen gat in het corpus.

\noindent De retrievallaag haalde 36 van de 140 aanwezige rapporten op, een recall van
25,7\,\%. De prestatie verschilt sterk per college, van 48,7\,\% voor de Onderwijsraad
tot 9,8\,\% voor de Adviesraad Internationale Vraagstukken
(Tabel~\ref{tab:dv1-advies-eval}). Dit cijfer is robuust. Toen alle zoekroutes
volledig werden aangezet, waaronder een trage inhoud-route, steeg het aantal
kandidaten van 1.135 naar 1.461, maar de recall bleef gelijk (24,3\,\%, binnen
hetzelfde betrouwbaarheidsinterval). Meer kandidaten leverden dus geen hogere recall op.

\noindent De oorzaak is structureel. Van de 140 aanwezige rapporten wordt de officiële
kopie maar 33 keer als kandidaat opgehaald, op welke positie in de lijst dan ook. De
overige 107, ofwel 76\,\%, komen nooit in de kandidaatlijst terecht. Dat ligt aan een
verschil tussen de zoekwijze en de manier waarop de rapporten zijn opgeslagen. De
routes zoeken op de collegenaam, terwijl de adviesrapporten in de officiële
bekendmakingen als bijlagen met hun eigen rapporttitel staan. Een datumvenster kan dit
niet verhelpen, want zo'n filter haalt kandidaten alleen weg en voegt er nooit toe.
Zoeken op de rapporttitel is evenmin een geldige oplossing, omdat de matcher juist nog
onbekende adviezen moet ontdekken en de titel dan niet bekend is.

\noindent De VLAM-poort is beoordeeld over de top-25 kandidaten per college, in totaal
125 documenten. Van de aanwezige rapporten die werden opgehaald en door VLAM
beoordeeld, kreeg elk het juiste label (VLAM-recall 100\,\%, 33 van de 33). De
precision van VLAM ligt op ten minste 30,2\,\%, gemeten tegen de website-set. Dat is
een ondergrens, want die set is onvolledig. Een handmatige controle van 15 van de 81
niet-gematchte positieven liet zien dat 7 echte adviesrapporten waren, afkomstig van de
Onderwijsraad en de Raad voor Cultuur en alleen niet aanwezig in de website-steekproef.
De overige 8 waren terechte valse positieven. Daarmee komt de precision na correctie op
ongeveer 63\,\%, tegenover 21,3\,\% voor de eerdere regelgebaseerde labeling. De valse
positieven concentreren zich bij de Adviesraad Internationale Vraagstukken en de Raad
voor het Openbaar Bestuur, juist de colleges met de laagste recall.

\noindent De recall van de retrievallaag vormt daarmee een bovengrens voor het hele
systeem. VLAM beoordeelt de opgehaalde kandidaten betrouwbaar, maar kan geen rapporten
labelen die de retrievallaag niet aanlevert. De uiteindelijke recall kan dus niet boven
de 25,7\,\% uitkomen die de retrievallaag haalt.

\noindent De schattingen zijn omgeven met onzekerheid. De steekproef is klein
(140 aanwezige rapporten, en 125 documenten voor de VLAM-beoordeling), en de
gecorrigeerde precision steunt op een handmatige herclassificatie van 15 items. De
grondwaarheid uit op websites gepubliceerde rapporten kan selectiebias bevatten. De
aanwezigheidsscan gebruikte bovendien de eerste 20.000 tekens van elk document, wat de
aanwezigheid van langere rapporten licht kan onderschatten.

\begin{table}[htbp]
\centering
\caption{Prestatie van de advies-matcher op een onafhankelijke grondwaarheid van
website-rapporten ($N = 140$ aantoonbaar aanwezige rapporten; de VLAM-poort werd over
125 documenten beoordeeld), met 95\,\%-betrouwbaarheidsintervallen volgens Wilson. Het
middelste blok geeft de retrieval-recall per college.}
\label{tab:dv1-advies-eval}
\footnotesize
\begin{tabular}{@{}lcc@{}}
\toprule
Metriek & Waarde & 95\,\%-BI \\
\midrule
Retrieval-recall & 25,7\,\% & 19,2--33,5 \\
Recall@5         & 4,3\,\%  & 2,0--9,0 \\
Recall@10        & 6,4\,\%  & 3,4--11,8 \\
Recall@25        & 14,3\,\% & 9,4--21,0 \\
\midrule
Onderwijsraad                                & 48,7\,\% & 33,9--63,8 \\
Raad voor Cultuur                            & 40,0\,\% & 19,8--64,3 \\
Raad voor het Openbaar Bestuur               & 18,8\,\% & 6,6--43,0 \\
Raad voor de Leefomgeving en Infrastructuur  & 15,8\,\% & 5,5--37,6 \\
Adviesraad Internationale Vraagstukken       & 9,8\,\%  & 4,3--21,0 \\
\midrule
VLAM-recall (gevoeligheid)                   & 100\,\%      & 89,6--100 \\
VLAM-precision (ondergrens)                  & 30,2\,\%     & 22,6--39,1 \\
VLAM-precision (na handmatige correctie)     & $\approx$\,63\,\% & --- \\
Regelgebaseerde precision (ter vergelijking) & 21,3\,\%     & --- \\
\bottomrule
\end{tabular}
\end{table}
